\section{Conclusion}\label{sec:conclusion}

This paper sets out, and partly executes, a programme for tracing physical flood risk from individual land parcels to bank capital in a setting where insurance is negligible and relief is the buffer. The institutional facts are established from public documents: a fixed calendar of official land valuations with no hazard rule, a collateral policy that does not reference those valuations, and relief rules that freeze the observed default of eligible borrowers for a known period. The framework makes explicit that the official valuer's capitalisation parameter and the bank's are different objects, that relief censors default and re-routes distress into restructuring and staging, and that the kink in loss given default makes unpriced hazard matter most for the safest-looking loans.

Empirically, the cross-section of \pilotParcelsM~million parcels against the 2011 flood shows a large urban value gap, most of which disappears when land cover measured ten years later is held fixed, leaving a residual of about 8\% that is imprecise and cannot be read as capitalisation; this is why the project does not rest on cross-sectional evidence. The pre-specified round-to-round design is complete up to the point where it needs data that do not yet exist: radar flood maps for two events, validated against optical imagery and against an independent analysis, hazard-comparable controls balanced within province with their residual gaps reported, minimum detectable effects of \regMdePooled{} log points in the primary sample and \evMdeW{} in the 2025 sample, and an estimator whose size is checked in \mcxReps{} Monte Carlo replications of both registered tests, a check that the wild cluster bootstrap passes and the placebo-footprint test fails until it is studentised, so that the plan is amended before registration to make the bootstrap primary and the placebo check studentised. The design is powered against realised-damage discounts and against a fully recognised hazard discount, not against partial recognition, and its estimate will be a first-round response rather than a capitalisation parameter. The credit and stress-test code recovers planted effects on a synthetic panel built to the requested schema, and the illustrative stress test shows that whether valuations price hazard moves the breaking point by 7--9\% at ordinary leverage, and by more at low leverage where the result depends on the assumed haircut and damage share.

The paper's content is therefore of three kinds and is labelled accordingly: established facts about institutions and data, a design with its power and its limits stated, and hypotheses about credit under relief that the design and the data request are built to test. The next valuation round, now likely in January 2028, will turn the first of those hypotheses into an estimate, provided its comparables postdate the floods.
